\uuid{YGdC}
\exo7id{3495}
\titre{exo7 3495}
\auteur{quercia}
\organisation{exo7}
\datecreate{2010-03-10}
\isIndication{false}
\isCorrection{true}
\chapitre{Application linéaire}
\sousChapitre{Morphismes particuliers}
\cpge{mpsi-mp2i.changements-bases.c, mpsi-mp2i.applications-lineaires.b, pcsi.applications-lineaires.b, ptsi.applications-lineaires.b, mpsi-mp2i.structures-algebriques.b, pc.complements-espaces-vectoriels.c, psi.complements-espaces-vectoriels.c, pt.complements-algebre-lineaire.c}
\module{Algèbre}
\niveau{L1}
\difficulte{}

\contenu{
\question{
Soit $E$ un $\R$-ev de dimension finie et $A$ une partie finie de $GL(E)$
stable par composition. On pose $u = \sum_{f\in A} f$. Montrer que
$\mathrm{tr}(u)=0  \Rightarrow  u=0$.
}
\reponse{
Si $A=\varnothing$ c'est évident.\par
Sinon, $A$ est un sous-groupe de $GL(E)$ donc $\frac{u}{\mathrm{Card}\, A}$ est
un projecteur et $\mathrm{tr}(u) = \mathrm{Card}\,(A)\mathrm{rg}(u)$.
}
}
